\subsection{严格态射与有限秩线性映射}

命题 1. --- 设 E 与 F 是局部凸空间，\(E_{1}\) 是 E 的有限余维闭子空间，而 u 是从 E 到 F 中的连续线性映射。要使 u 是具有闭像的严格态射，必须且只需 \(u|E_{1}\) 是这样的态射。

首先假设线性映射 \(u|E_1\) 是单射。于是我们有 \(E_1 \cap \text{Ker}(u) = \{0\}\)，从而 \(\text{Ker}(u)\) 是有限维的。设 S 是 E 的一个向量子空间，它是 \(E_1 + \text{Ker}(u)\) 的补。空间 E 是 \(E_1\)、\(\text{Ker}(u)\) 与 S 的拓扑直和（EVT, I, p. 15, cor. 4 与 prop. 3）。如果 \(u\) 是具有闭像的严格态射，那么它通过限制定义了从 \(E_1 \oplus S\) 到 F 的一个闭向量子空间上的同构，从而更加必然地定义了从 \(E_1\) 到 F 的一个闭向量子空间上的同构。反之，假设 \(u\) 通过限制定义了从 \(E_1\) 到 \(u(E_1)\) 上的同构，且 \(u(E_1)\) 在 F 中是闭的。我们有 \(u(E) = u(E_1) \oplus u(S)\)。由此得出 \(u(E)\) 是 \(u(E_1)\) 与 \(u(S)\) 的拓扑直和，且在 F 中是闭的（loc. cit.）。由于 \(u(E_1)\) 是闭的，且我们有 \(u(E_1) \cap u(S) = \{0\}\)，空间 \(u(S)\) 是分离的，且 \(u\) 通过限制定义了从 S 到 \(u(S)\) 上的同构（EVT, I, p. 15, prop. 3 的 cor.），因而也定义了从 \(E_1 \oplus S\) 到 \(u(E)\) 上的同构。这就证明了 \(u\) 是具有闭像的严格态射。

转到一般情形，置 \(N = E_{1} \cap \operatorname{Ker}(u)\) 与 G = E/N。局部凸空间 \(E_{1}/N\) 等同于 G 的一个有限余维闭向量子空间 \(G_{1}\)。用 \(v\colon G \to F\) 表示由 u 通过过渡到商而得到的连续线性映射。要使 \(u\)（相应地，\(u|E_{1}\)）是具有闭像的严格态射，必须且只需 \(v\)（相应地，\(v|G_{1}\)）是这样的态射。这就化归到已处理过的情形。

推论 1. --- 假设 F 是分离的。设 \(v \in \mathcal{L}^{\mathrm{f}}(\mathrm{E};\mathrm{F})\)。如果 u 是从 E 到 F 中的具有闭像的严格态射，那么 \(u + v\) 也是。

由于 F 是分离的，v 的核是 E 的闭向量子空间；它在 E 中有有限余维，且 \(u + v\) 与 \(u\) 在 \(\operatorname{Ker}(v)\) 上有相同的限制。推论于是由该命题得出。

推论 2. --- 设 T 是 F 的分离的有限维向量子空间。用 \(\pi\colon\mathrm{F}\to\mathrm{F}/\mathrm{T}\) 表示典范满射。要使 u 是具有闭像的严格态射，必须且只需 \(\pi\circ u\) 是这样的态射。

T 到自身的恒等映射延拓为连续映射 \(q \colon \mathrm{F} \to \mathrm{T}\)（EVT, II, p. 26, remark）。\(q\) 的核 S 是 T 的闭拓扑补。设 \(p \colon \mathrm{F} \to \mathrm{F}\) 是与分解 \(\mathrm{F} = \mathrm{T} \oplus \mathrm{S}\) 相伴的、以 S 为像的投影。要使 \(\pi \circ u\) 是具有闭像的严格态射，必须且只需 \(p \circ u\) 是这样的态射。此外，\(p \circ u\) 与 \(u\) 在 \(u^{-1}(\mathrm{S})\)（它是 E 中的有限余维闭向量子空间）上有相同的限制，断言于是由该命题得出。

